% Strona tytułowa sprawozdania według wzoru SAN (Word): logo, kierunek, dane
% studenta po prawej, przedmiot, rodzaj pracy, tytuł, prowadzący, miasto i rok.
%
% Użycie w main.tex (po \documentclass{san}, przed \begin{document}):
%   % Strona tytułowa sprawozdania według wzoru SAN (Word): logo, kierunek, dane
% studenta po prawej, przedmiot, rodzaj pracy, tytuł, prowadzący, miasto i rok.
%
% Użycie w main.tex (po \documentclass{san}, przed \begin{document}):
%   % Strona tytułowa sprawozdania według wzoru SAN (Word): logo, kierunek, dane
% studenta po prawej, przedmiot, rodzaj pracy, tytuł, prowadzący, miasto i rok.
%
% Użycie w main.tex (po \documentclass{san}, przed \begin{document}):
%   % Strona tytułowa sprawozdania według wzoru SAN (Word): logo, kierunek, dane
% studenta po prawej, przedmiot, rodzaj pracy, tytuł, prowadzący, miasto i rok.
%
% Użycie w main.tex (po \documentclass{san}, przed \begin{document}):
%   \input{strona-tytulowa}
%   + dane: \przedmiot, \tytul, \nralbumu, \grupa, \prowadzacy, \miasto, \rok (z san.cls)
%     oraz poniższe (mają wartości domyślne, zmieniaj tylko potrzebne):
%   \kierunek{Informatyka}
%   \nazwiskoimie{Roman Piotr}
%   \grupaslownie{IV}            – pusta wartość {} ukrywa wiersz
%   \trybstudiow{Studia Inż. Zaoczne}   – pusta wartość {} ukrywa wiersz
%   \rodzajpracy{Sprawozdanie}
%   \zajecia{z projektu}         – np. {z Laboratorium nr 1}; pusta ukrywa wiersz
%   \logo{san-logo.jpg}          – plik w img/
% Strona powstaje przez \maketitlepage (zastępuje wersję z san.cls, pliku klasy nie zmienia).

\makeatletter

\newcommand{\kierunek}[1]{\gdef\@kierunek{#1}}
\newcommand{\nazwiskoimie}[1]{\gdef\@nazwiskoimie{#1}}
\newcommand{\grupaslownie}[1]{\gdef\@grupaslownie{#1}}
\newcommand{\trybstudiow}[1]{\gdef\@trybstudiow{#1}}
\newcommand{\rodzajpracy}[1]{\gdef\@rodzajpracy{#1}}
\newcommand{\zajecia}[1]{\gdef\@zajecia{#1}}
\newcommand{\logo}[1]{\gdef\@logo{#1}}

\kierunek{Informatyka}
\nazwiskoimie{Roman Piotr}
\grupaslownie{IV}
\trybstudiow{Studia Inż. Zaoczne}
\rodzajpracy{Sprawozdanie}
\zajecia{z projektu}
\logo{san-logo.jpg}

% Wiersz danych studenta; pomijany, gdy wartość jest pusta.
\newcommand{\@tpwiersz}[1]{%
  \if\relax\detokenize{#1}\relax\else #1\par\vspace{0.15cm}\fi}

\renewcommand{\maketitlepage}{%
  \begin{titlepage}
    \thispagestyle{empty}
    \centering
    \includegraphics[width=0.72\textwidth]{\@logo}\par
    \vspace{1.2cm}
    {\LARGE Kierunek studiów: \textbf{\@kierunek}\par}
    \vspace{0.5cm}
    \hspace*{0.3\textwidth}%
    \begin{minipage}{0.6\textwidth}
      \raggedright\small
      \@tpwiersz{Nazwisko i Imię: \@nazwiskoimie}
      \@tpwiersz{Numer albumu: \@nralbumu}
      \if\relax\detokenize\expandafter{\@grupaslownie}\relax\else
        \@tpwiersz{Grupa: \@grupaslownie}\fi
      \@tpwiersz{nr grupy: \@grupa}
      \if\relax\detokenize\expandafter{\@trybstudiow}\relax\else
        \@tpwiersz{\@trybstudiow}\fi
    \end{minipage}\par
    \vspace{1.3cm}
    {\sffamily\bfseries\Large \@przedmiot\par}
    \vspace{0.7cm}
    {\sffamily\bfseries\Large \@rodzajpracy\par}
    \if\relax\detokenize\expandafter{\@zajecia}\relax\else
      \vspace{0.5cm}
      {\sffamily\bfseries\Large \@zajecia\par}
    \fi
    \vspace{0.8cm}
    {\bfseries\LARGE \@tytul\par}
    \vfill
    Prowadzący zajęcia: \textbf{\@prowadzacy}\par
    \vspace{2cm}
    {\bfseries\Large \@miasto\ \@rok\par}
    \vspace{1cm}
  \end{titlepage}%
}

\makeatother

%   + dane: \przedmiot, \tytul, \nralbumu, \grupa, \prowadzacy, \miasto, \rok (z san.cls)
%     oraz poniższe (mają wartości domyślne, zmieniaj tylko potrzebne):
%   \kierunek{Informatyka}
%   \nazwiskoimie{Roman Piotr}
%   \grupaslownie{IV}            – pusta wartość {} ukrywa wiersz
%   \trybstudiow{Studia Inż. Zaoczne}   – pusta wartość {} ukrywa wiersz
%   \rodzajpracy{Sprawozdanie}
%   \zajecia{z projektu}         – np. {z Laboratorium nr 1}; pusta ukrywa wiersz
%   \logo{san-logo.jpg}          – plik w img/
% Strona powstaje przez \maketitlepage (zastępuje wersję z san.cls, pliku klasy nie zmienia).

\makeatletter

\newcommand{\kierunek}[1]{\gdef\@kierunek{#1}}
\newcommand{\nazwiskoimie}[1]{\gdef\@nazwiskoimie{#1}}
\newcommand{\grupaslownie}[1]{\gdef\@grupaslownie{#1}}
\newcommand{\trybstudiow}[1]{\gdef\@trybstudiow{#1}}
\newcommand{\rodzajpracy}[1]{\gdef\@rodzajpracy{#1}}
\newcommand{\zajecia}[1]{\gdef\@zajecia{#1}}
\newcommand{\logo}[1]{\gdef\@logo{#1}}

\kierunek{Informatyka}
\nazwiskoimie{Roman Piotr}
\grupaslownie{IV}
\trybstudiow{Studia Inż. Zaoczne}
\rodzajpracy{Sprawozdanie}
\zajecia{z projektu}
\logo{san-logo.jpg}

% Wiersz danych studenta; pomijany, gdy wartość jest pusta.
\newcommand{\@tpwiersz}[1]{%
  \if\relax\detokenize{#1}\relax\else #1\par\vspace{0.15cm}\fi}

\renewcommand{\maketitlepage}{%
  \begin{titlepage}
    \thispagestyle{empty}
    \centering
    \includegraphics[width=0.72\textwidth]{\@logo}\par
    \vspace{1.2cm}
    {\LARGE Kierunek studiów: \textbf{\@kierunek}\par}
    \vspace{0.5cm}
    \hspace*{0.3\textwidth}%
    \begin{minipage}{0.6\textwidth}
      \raggedright\small
      \@tpwiersz{Nazwisko i Imię: \@nazwiskoimie}
      \@tpwiersz{Numer albumu: \@nralbumu}
      \if\relax\detokenize\expandafter{\@grupaslownie}\relax\else
        \@tpwiersz{Grupa: \@grupaslownie}\fi
      \@tpwiersz{nr grupy: \@grupa}
      \if\relax\detokenize\expandafter{\@trybstudiow}\relax\else
        \@tpwiersz{\@trybstudiow}\fi
    \end{minipage}\par
    \vspace{1.3cm}
    {\sffamily\bfseries\Large \@przedmiot\par}
    \vspace{0.7cm}
    {\sffamily\bfseries\Large \@rodzajpracy\par}
    \if\relax\detokenize\expandafter{\@zajecia}\relax\else
      \vspace{0.5cm}
      {\sffamily\bfseries\Large \@zajecia\par}
    \fi
    \vspace{0.8cm}
    {\bfseries\LARGE \@tytul\par}
    \vfill
    Prowadzący zajęcia: \textbf{\@prowadzacy}\par
    \vspace{2cm}
    {\bfseries\Large \@miasto\ \@rok\par}
    \vspace{1cm}
  \end{titlepage}%
}

\makeatother

%   + dane: \przedmiot, \tytul, \nralbumu, \grupa, \prowadzacy, \miasto, \rok (z san.cls)
%     oraz poniższe (mają wartości domyślne, zmieniaj tylko potrzebne):
%   \kierunek{Informatyka}
%   \nazwiskoimie{Roman Piotr}
%   \grupaslownie{IV}            – pusta wartość {} ukrywa wiersz
%   \trybstudiow{Studia Inż. Zaoczne}   – pusta wartość {} ukrywa wiersz
%   \rodzajpracy{Sprawozdanie}
%   \zajecia{z projektu}         – np. {z Laboratorium nr 1}; pusta ukrywa wiersz
%   \logo{san-logo.jpg}          – plik w img/
% Strona powstaje przez \maketitlepage (zastępuje wersję z san.cls, pliku klasy nie zmienia).

\makeatletter

\newcommand{\kierunek}[1]{\gdef\@kierunek{#1}}
\newcommand{\nazwiskoimie}[1]{\gdef\@nazwiskoimie{#1}}
\newcommand{\grupaslownie}[1]{\gdef\@grupaslownie{#1}}
\newcommand{\trybstudiow}[1]{\gdef\@trybstudiow{#1}}
\newcommand{\rodzajpracy}[1]{\gdef\@rodzajpracy{#1}}
\newcommand{\zajecia}[1]{\gdef\@zajecia{#1}}
\newcommand{\logo}[1]{\gdef\@logo{#1}}

\kierunek{Informatyka}
\nazwiskoimie{Roman Piotr}
\grupaslownie{IV}
\trybstudiow{Studia Inż. Zaoczne}
\rodzajpracy{Sprawozdanie}
\zajecia{z projektu}
\logo{san-logo.jpg}

% Wiersz danych studenta; pomijany, gdy wartość jest pusta.
\newcommand{\@tpwiersz}[1]{%
  \if\relax\detokenize{#1}\relax\else #1\par\vspace{0.15cm}\fi}

\renewcommand{\maketitlepage}{%
  \begin{titlepage}
    \thispagestyle{empty}
    \centering
    \includegraphics[width=0.72\textwidth]{\@logo}\par
    \vspace{1.2cm}
    {\LARGE Kierunek studiów: \textbf{\@kierunek}\par}
    \vspace{0.5cm}
    \hspace*{0.3\textwidth}%
    \begin{minipage}{0.6\textwidth}
      \raggedright\small
      \@tpwiersz{Nazwisko i Imię: \@nazwiskoimie}
      \@tpwiersz{Numer albumu: \@nralbumu}
      \if\relax\detokenize\expandafter{\@grupaslownie}\relax\else
        \@tpwiersz{Grupa: \@grupaslownie}\fi
      \@tpwiersz{nr grupy: \@grupa}
      \if\relax\detokenize\expandafter{\@trybstudiow}\relax\else
        \@tpwiersz{\@trybstudiow}\fi
    \end{minipage}\par
    \vspace{1.3cm}
    {\sffamily\bfseries\Large \@przedmiot\par}
    \vspace{0.7cm}
    {\sffamily\bfseries\Large \@rodzajpracy\par}
    \if\relax\detokenize\expandafter{\@zajecia}\relax\else
      \vspace{0.5cm}
      {\sffamily\bfseries\Large \@zajecia\par}
    \fi
    \vspace{0.8cm}
    {\bfseries\LARGE \@tytul\par}
    \vfill
    Prowadzący zajęcia: \textbf{\@prowadzacy}\par
    \vspace{2cm}
    {\bfseries\Large \@miasto\ \@rok\par}
    \vspace{1cm}
  \end{titlepage}%
}

\makeatother

%   + dane: \przedmiot, \tytul, \nralbumu, \grupa, \prowadzacy, \miasto, \rok (z san.cls)
%     oraz poniższe (mają wartości domyślne, zmieniaj tylko potrzebne):
%   \kierunek{Informatyka}
%   \nazwiskoimie{Roman Piotr}
%   \grupaslownie{IV}            – pusta wartość {} ukrywa wiersz
%   \trybstudiow{Studia Inż. Zaoczne}   – pusta wartość {} ukrywa wiersz
%   \rodzajpracy{Sprawozdanie}
%   \zajecia{z projektu}         – np. {z Laboratorium nr 1}; pusta ukrywa wiersz
%   \logo{san-logo.jpg}          – plik w img/
% Strona powstaje przez \maketitlepage (zastępuje wersję z san.cls, pliku klasy nie zmienia).

\makeatletter

\newcommand{\kierunek}[1]{\gdef\@kierunek{#1}}
\newcommand{\nazwiskoimie}[1]{\gdef\@nazwiskoimie{#1}}
\newcommand{\grupaslownie}[1]{\gdef\@grupaslownie{#1}}
\newcommand{\trybstudiow}[1]{\gdef\@trybstudiow{#1}}
\newcommand{\rodzajpracy}[1]{\gdef\@rodzajpracy{#1}}
\newcommand{\zajecia}[1]{\gdef\@zajecia{#1}}
\newcommand{\logo}[1]{\gdef\@logo{#1}}

\kierunek{Informatyka}
\nazwiskoimie{Roman Piotr}
\grupaslownie{IV}
\trybstudiow{Studia Inż. Zaoczne}
\rodzajpracy{Sprawozdanie}
\zajecia{z projektu}
\logo{san-logo.jpg}

% Wiersz danych studenta; pomijany, gdy wartość jest pusta.
\newcommand{\@tpwiersz}[1]{%
  \if\relax\detokenize{#1}\relax\else #1\par\vspace{0.15cm}\fi}

\renewcommand{\maketitlepage}{%
  \begin{titlepage}
    \thispagestyle{empty}
    \centering
    \includegraphics[width=0.72\textwidth]{\@logo}\par
    \vspace{1.2cm}
    {\LARGE Kierunek studiów: \textbf{\@kierunek}\par}
    \vspace{0.5cm}
    \hspace*{0.3\textwidth}%
    \begin{minipage}{0.6\textwidth}
      \raggedright\small
      \@tpwiersz{Nazwisko i Imię: \@nazwiskoimie}
      \@tpwiersz{Numer albumu: \@nralbumu}
      \if\relax\detokenize\expandafter{\@grupaslownie}\relax\else
        \@tpwiersz{Grupa: \@grupaslownie}\fi
      \@tpwiersz{nr grupy: \@grupa}
      \if\relax\detokenize\expandafter{\@trybstudiow}\relax\else
        \@tpwiersz{\@trybstudiow}\fi
    \end{minipage}\par
    \vspace{1.3cm}
    {\sffamily\bfseries\Large \@przedmiot\par}
    \vspace{0.7cm}
    {\sffamily\bfseries\Large \@rodzajpracy\par}
    \if\relax\detokenize\expandafter{\@zajecia}\relax\else
      \vspace{0.5cm}
      {\sffamily\bfseries\Large \@zajecia\par}
    \fi
    \vspace{0.8cm}
    {\bfseries\LARGE \@tytul\par}
    \vfill
    Prowadzący zajęcia: \textbf{\@prowadzacy}\par
    \vspace{2cm}
    {\bfseries\Large \@miasto\ \@rok\par}
    \vspace{1cm}
  \end{titlepage}%
}

\makeatother
